%% ===================================================================
\section{Light-cone wave functions to order \texorpdfstring{$g^3$}{g3}}\label{sec:lcwf}
%% ===================================================================

In this section we list the LCPT wave functions from which the coefficients of $\Omega$ are obtained. All results are given in transverse coordinate space; the Fourier transforms used are listed in App.~\ref{app:conv}, Eqs.~\eqref{eq:i1}--\eqref{eq:i9}. Plus momenta are bounded by the soft window: every soft gluon has $\Lambda<p^+<\vee$. For a pair produced from a parent with momentum $k^+$ this restricts the momentum fraction to $\Lambda/k^+<\xi<1-\Lambda/k^+$. The transverse dimension is $d_\perp=2-\epsilon$ wherever a loop integral requires regularization, and $\bar\mu^2$ is the $\MSbar$ scale.

%% -------------------------------------------------------------------
\subsection{The valence vacuum}\label{sec:vac}

To order $g^3$ the dressed valence state is
\begin{equation}
|\Psi\rangle=\mf N|0\rangle+|\Psi^{\rm LO}_{g\rho}\rangle+|\Psi_{gg\rho}\rangle+|\Psi_{gg\rho\rho}\rangle+|\Psi_{g\rho}\rangle+|\Psi_{g\rho\rho}\rangle+|\Psi_{g\rho\rho\rho}\rangle,
\label{eq:NLOWF}
\end{equation}
where the subscript counts the soft gluons ($g$) and the valence charges ($\rho$) of each component.

\paragraph*{One gluon at order $g$.} A valence charge emits one soft gluon:
\begin{equation}
|\Psi^{\rm LO}_{g\rho}\rangle=\frac{ig}{2\pi^{3/2}}\int_{\Lambda}^{\vee}\frac{dk^+}{\sqrt{k^+}}\int_{x,z}\KK^i(x-z)\,\rho^a(x)\,|g^a_i(k^+,z)\rangle .
\label{eq:LOwf}
\end{equation}

\paragraph*{Two gluons, one charge.} Three processes contribute at order $g^2$. A gluon emitted by the charge splits into two through the three-gluon vertex; the charge emits two gluons through the instantaneous interaction; and the commutator part of two successive emissions from the same charge has the same colour structure. Their sum is
\begin{align}
|\Psi_{gg\rho}\rangle={}&\int_{\Lambda}^{\vee}dk^+\int_{\Lambda/k^+}^{1-\Lambda/k^+}d\xi\int_{x,z,w}\frac{ig^2f^{abc}\rho^c(x)\sqrt{\xi(1-\xi)}}{8\pi^3\big[(1-\xi)(x-w)^2+\xi(x-z)^2\big]}\notag\\
&\times\Bigg[\frac{\delta^{ij}\big[(x-z)^2-(x-w)^2\big]}{2(z-w)^2}+\frac{(x-w)^i(z-w)^j}{\xi\,(z-w)^2}+\frac{(x-z)^j(z-w)^i}{(1-\xi)(z-w)^2}\notag\\
&\hspace{2.5em}+\frac{(x-z)^j(x-w)^i}{2\xi\,(x-z)^2}-\frac{(x-z)^j(x-w)^i}{2(1-\xi)(x-w)^2}\Bigg]\,|g^b_j(\xi k^+,z)\,g^a_i((1-\xi)k^+,w)\rangle .
\label{eq:ggrho}
\end{align}
For $\xi\to0$ the gluon at $z$ is soft, and the bracket divided by $(1-\xi)(x-w)^2+\xi(x-z)^2$ reduces to $\xi^{-1}\KK^i(x-w)\big[-\KK^j(w-z)+\tfrac12\KK^j(x-z)\big]$. The first term is the soft gluon emitted by the harder gluon at $w$. The second is one half of the commutator of two independent emissions, $\rho^b(x)\rho^c(y)=\frac12\{\rho^b(x),\rho^c(y)\}+\frac12[\rho^b(x),\rho^c(y)]$.

\paragraph*{Two gluons, two charges.} The symmetric (anticommutator) part of two independent emissions is
\begin{align}
|\Psi_{gg\rho\rho}\rangle={}&-\int_{\Lambda}^{\vee}dk^+\int_{\Lambda/k^+}^{1-\Lambda/k^+}d\xi\int_{x,y,z,w}\frac{g^2\{\rho^a(y),\rho^b(x)\}\,(y-z)^i(x-w)^j}{16\pi^3\sqrt{\xi(1-\xi)}\,(y-z)^2(x-w)^2}\notag\\
&\times|g^b_j((1-\xi)k^+,w)\,g^a_i(\xi k^+,z)\rangle .
\label{eq:ggrhorho}
\end{align}

\paragraph*{One gluon, one charge, at order $g^3$.} This is the one-loop correction to the emission~\eqref{eq:LOwf}. It contains the self-energy of the emitted gluon (gluon loop), the vertex corrections in which the emitted gluon splits and recombines through the instantaneous interaction, and the corresponding diagrams with the loop attached to the valence charge. After the $\xi$ integration of the loop,
\begin{align}
|\Psi_{g\rho}\rangle={}&-\int_{\Lambda}^{\vee}\frac{dk^+}{\sqrt{k^+}}\int_{x,z}\frac{ig^3N_c\,\rho^a(x)\,(x-z)^i}{32\pi^{7/2}(x-z)^2}\Bigg\{\Big[-\frac2\epsilon-2\gamma_E-\log\frac{(x-z)^2\bar\mu^2}{4}\Big]\notag\\
&\hspace{6em}\times\Big[\frac{11}{3}+\log\frac{\Lambda}{k^+}+\log\Big(\frac{\vee}{k^+}-1\Big)-\log\frac{\vee}{\Lambda}\Big]\notag\\
&\hspace{6em}+\frac{\pi^2}{3}-\frac{67}{9}+\log^2\Big(1-\frac{k^+}{\vee}\Big)+\log\Big(1-\frac{k^+}{\vee}\Big)\log\frac{\vee}{k^+}\Bigg\}|g^a_i(k^+,z)\rangle .
\label{eq:grho}
\end{align}
Three features of Eq.~\eqref{eq:grho} recur throughout the paper.
\begin{itemize}
\item The UV pole multiplies $\frac{11}{3}=\beta_0/N_c$ (for $n_f=0$). The $\xi$-integrand of the gluon loop is the real gluon splitting function, see Sec.~\ref{sec:rc}.
\item The rapidity logarithms combine as
\begin{equation}
\log\frac{\Lambda}{k^+}+\log\Big(\frac{\vee}{k^+}-1\Big)-\log\frac{\vee}{\Lambda}=-2\log\frac{k^+}{\Lambda}+\log\Big(1-\frac{k^+}{\vee}\Big),
\label{eq:grhologs}
\end{equation}
so the loop contains only the logarithm of the region below $k^+$, plus a term that vanishes as $k^+/\vee\to0$.
\item The constant $\frac{\pi^2}{3}-\frac{67}{9}=-2\big(\frac{67}{18}-\frac{\pi^2}{6}\big)=-2K/C_A$ is the Catani--Marchesini--Webber constant~\cite{Catani:1990rr} (for $n_f=0$).
\end{itemize}

\paragraph*{One gluon, two charges.} In a mixed representation, with the Fourier variables $\mathbf k$ (conjugate to $z-y$) and $\mathbf p$ (conjugate to $y-x$),
\begin{align}
|\Psi_{g\rho\rho}\rangle={}&\frac{1}{2\pi}\int_{\Lambda}^{\vee}\frac{dk^+}{\sqrt{k^+}}\int d^2\mathbf k\,d^2\mathbf p\int_{z,x,y}e^{i\mathbf k\cdot(z-y)}e^{i\mathbf p\cdot(y-x)}\,\frac{ig^3f^{abc}\{\rho^b(x),\rho^c(y)\}}{32\pi^{9/2}\,\mathbf p^2\mathbf k^2}\notag\\
&\times\Bigg[\log\frac{\Lambda}{k^+}\,\frac{\mathbf k^2\mathbf p^i-(\mathbf k\cdot\mathbf p)\,\mathbf k^i}{(\mathbf k-\mathbf p)^2}-2\mathbf k^i\log\frac{k^+\mathbf p^2}{\vee\,\mathbf k^2+k^+\mathbf p^2}-2\mathbf k^i\log\frac{\vee}{\Lambda}\notag\\
&\qquad+\frac{(2\mathbf p^2-2\mathbf k\cdot\mathbf p+\mathbf k^2)\,\mathbf k^i+\mathbf k^2\mathbf p^i}{(\mathbf k-\mathbf p)^2}\log\Big[1+\Big(\frac{\vee}{k^+}-1\Big)\frac{\mathbf k^2}{(\mathbf k-\mathbf p)^2}\Big]\notag\\
&\qquad-\frac{\mathbf p^2\mathbf k^i+\mathbf k^2\mathbf p^i}{(\mathbf k-\mathbf p)^2}\log\frac{\vee-k^+}{\Lambda}-\mathbf k^i\log\frac{\vee}{k^+}\Bigg]|g^a_i(k^+,z)\rangle ,
\label{eq:grhorho}
\end{align}
where terms that vanish for $\Lambda\to0$ at fixed $k^+$ have been dropped.

\paragraph*{One gluon, three charges.} The product of the normalization correction of one charge pair with the emission by a third charge gives
\begin{equation}
|\Psi_{g\rho\rho\rho}\rangle=-\log\frac{\vee}{\Lambda}\int_{\Lambda}^{\vee}\frac{dk^+}{\sqrt{k^+}}\int_{w,x,y,u,z}\frac{ig^3\,\rho^b(x)\rho^b(y)\rho^a(w)\,(x-u)\cdot(y-u)\,(w-z)^i}{16\pi^{9/2}(x-u)^2(y-u)^2(w-z)^2}\,|g^a_i(k^+,z)\rangle .
\label{eq:grhorhorho}
\end{equation}

\paragraph*{Normalization.} To order $g^2$,
\begin{equation}
\mf N=1-\frac{g^2}{8\pi^3}\log\frac{\vee}{\Lambda}\int_{x,y,z}K(x,y,z)\,\rho^a(x)\rho^a(y),\qquad K(x,y,z)\equiv\KK(x-z)\cdot\KK(y-z)=\frac{(x-z)\cdot(y-z)}{(x-z)^2(y-z)^2}.
\label{eq:NX}
\end{equation}
The kernel $K(x,y,z)$ is symmetric in its first two arguments. It is the BFKL/JIMWLK kernel used throughout the paper.

%% -------------------------------------------------------------------
\subsection{A single soft gluon}\label{sec:singleglue}

We now dress a soft gluon $|g^a_i(k^+,z')\rangle$ to order $g^2$. We write $X'\equiv x-z'$ for the distance of a valence charge at $x$ from the incoming gluon, and $X\equiv x-z$ for its distance from the outgoing gluon at $z$. The outgoing momentum fraction is $\xi=p^+/k^+$. Contributions in which the outgoing gluon coincides with the incoming one are part of $\mf N$ and are not included below.

\paragraph*{No gluon in the final state.} The valence charge absorbs the gluon:
\begin{equation}
|\Psi^a_i(k^+,z')\rangle_\rho=\frac{ig}{2\pi^{3/2}\sqrt{k^+}}\int_x\KK^i(X')\,\rho^a(x)\,|0\rangle .
\label{eq:absorb}
\end{equation}

\paragraph*{One gluon in the final state: one charge.} The soft gluon exchanges plus momentum with the valence charge. The outgoing gluon can be harder ($\xi>1$, momentum flows from the charge to the gluon) or softer ($\xi<1$). The instantaneous interaction~\eqref{eq:m5} gives the same integrand in both cases:
\begin{equation}
|\Psi^a_i(k^+,z')\rangle^{\rm inst}_{g\rho}=\frac{ig^2f^{abc}}{(2\pi)^3}\bigg[\int_{1+\Lambda/k^+}^{\vee/k^+}+\int_{\Lambda/k^+}^{1-\Lambda/k^+}\bigg]d\xi\int_{x,z}\frac{\sqrt\xi(1+\xi)}{(1-\xi)^2}\,\frac{\rho^c(x)}{\xi X^2-X'^2}\,|g^b_i(\xi k^+,z)\rangle .
\label{eq:inst}
\end{equation}
The non-instantaneous contributions, in which the charge emits (absorbs) a gluon that then merges with (splits off) the incoming gluon through the three-gluon vertex, are
\begin{align}
|\Psi^a_i(k^+,z')\rangle^{(1)}_{g\rho}={}&-\frac{ig^2f^{abc}}{4\pi^3}\int_{1+\Lambda/k^+}^{\vee/k^+}d\xi\int_{x,z}\frac{\sqrt\xi}{1-\xi}\,\frac{X'^m-\xi X^m}{X'^2-\xi X^2}\,\frac{X'^l-X^l}{(X'-X)^2}\,\rho^c(x)\notag\\
&\hspace{6em}\times\Big(\delta_{il}\delta_{jm}-\frac1\xi\delta_{jl}\delta_{im}-\frac{1}{1-\xi}\delta_{ij}\delta_{lm}\Big)|g^b_j(\xi k^+,z)\rangle\qquad(\xi>1),\label{eq:nonInst1}\\
|\Psi^a_i(k^+,z')\rangle^{(2)}_{g\rho}={}&-\frac{ig^2f^{abc}}{4\pi^3}\int_{\Lambda/k^+}^{1-\Lambda/k^+}d\xi\int_{x,z}\frac{\xi^{3/2}}{1-\xi}\,\frac{X'^l-X^l}{X'^2-\xi X^2}\,\frac{X'^m-\xi X^m}{(X'-\xi X)^2}\,\rho^c(x)\notag\\
&\hspace{6em}\times\Big(\delta_{il}\delta_{jm}-\frac1\xi\delta_{jl}\delta_{im}-\frac{1}{1-\xi}\delta_{lm}\delta_{ij}\Big)|g^b_j(\xi k^+,z)\rangle\qquad(\xi<1).\label{eq:nonInst2}
\end{align}

\paragraph*{The endpoint $\xi\to1$.} All three expressions contain $(1-\xi)^{-2}$. The integration limits exclude only the window $|\xi-1|<\Lambda/k^+$, so a surviving double pole would give a power-like term $\propto k^+/\Lambda$. Near $\xi=1$ the instantaneous term and the $\delta_{ij}\delta_{lm}/(1-\xi)$ parts of Eqs.~\eqref{eq:nonInst1} and~\eqref{eq:nonInst2} behave as
\begin{equation}
\frac{ig^2f^{abc}}{4\pi^3}\,\frac{\delta_{ij}\,\rho^c(x)}{(1-\xi)^2\,(X^2-X'^2)}\times
\begin{cases}+1&\text{(instantaneous, both sides)},\\ -1&\text{[Eq.~\eqref{eq:nonInst1}, $\xi\to1^+$, and Eq.~\eqref{eq:nonInst2}, $\xi\to1^-$]},\end{cases}
\label{eq:endpoint}
\end{equation}
so the double pole cancels on both sides of $\xi=1$. This fixes the overall sign of Eq.~\eqref{eq:nonInst1}. It can also be checked directly from the matrix elements~\eqref{eq:m2} and~\eqref{eq:m3}. For $\xi>1$ the soft gluon of momentum $p^+-k^+$ is exchanged between the valence charge and the hard gluon. Its eikonal coupling to the hard gluon, the term $\frac{p^++k^+}{p^+-k^+}(p-k)^m\delta_{ij}$ of the three-gluon vertex, combined with the exact energy denominator $E_k-E_{k,p-k}=-(\mathbf p-\mathbf k)^2/2(p^+-k^+)$, reproduces the instantaneous term~\eqref{eq:inst} exactly, with the opposite sign. With the opposite overall sign of Eq.~\eqref{eq:nonInst1}, the instantaneous and non-instantaneous double poles would add for $\xi>1$. In the cross section this would leave a term $\propto1/\Lambda$ (Sec.~\ref{sec:finite}).
\note{The sign of Eq.~\eqref{eq:nonInst1} is opposite to the one in the draft notes (Eq.~onegrho12 and its coordinate form). The draft's momentum-space form in the variable $\xi$ already has the sign used here; onegrho12 and the coordinate form should be re-derived to confirm.}
%TODO: re-derive onegrho12 (p+>k+ non-instantaneous) and confirm the overall sign; propagate to B3 (done here).

\paragraph*{One gluon in the final state: two charges.} The incoming gluon is absorbed by one charge and a new gluon is emitted by another, in either order. Splitting the product of charges into its symmetric and antisymmetric parts,
\begin{align}
|\Psi^a_i(k^+,z')\rangle_{g\{\rho,\rho\}}&=-\frac{g^2}{8\pi^3}\int_{\Lambda/k^+}^{\vee/k^+}\frac{d\xi}{\sqrt\xi}\int_{x,y,z}\frac{X'^iY^j}{X'^2Y^2}\,\{\rho^a(x),\rho^b(y)\}\,|g^b_j(\xi k^+,z)\rangle,\qquad Y\equiv y-z,\label{eq:g2rhoS}\\
|\Psi^a_i(k^+,z')\rangle_{g[\rho,\rho]}&=-\frac{ig^2f^{abc}}{8\pi^3}\int_{\Lambda/k^+}^{\vee/k^+}\frac{d\xi}{\sqrt\xi}\int_{x,z}\frac{X'^iX^j}{X'^2X^2}\,\frac{X'^2+\xi X^2}{\xi X^2-X'^2}\,\rho^c(x)\,|g^b_j(\xi k^+,z)\rangle .\label{eq:g2rhoA}
\end{align}

\paragraph*{Two gluons in the final state.} The incoming gluon splits through the three-gluon vertex, independently of the valence charges:
\begin{align}
|\Psi^a_i(k^+,z')\rangle_{gg}={}&-\frac{gf^{abc}\sqrt{k^+}}{4\pi^{3/2}}\int_z\int_{\Lambda/k^+}^{1-\Lambda/k^+}d\xi\,\sqrt{\frac{1-\xi}{\xi}}\,\frac{(z-z')^m}{(z-z')^2}\Big[\delta_{im}\delta_{jl}-\frac1\xi\delta_{jm}\delta_{il}-\frac{1}{1-\xi}\delta_{lm}\delta_{ij}\Big]\notag\\
&\times|g^c_l((1-\xi)k^+,z)\,g^b_j\big(\xi k^+,\tfrac{z'-(1-\xi)z}{\xi}\big)\rangle .
\label{eq:gg}
\end{align}
The transverse centre of plus momentum of the pair is at the position $z'$ of the parent. A valence charge can also emit a second gluon independently of the incoming one:
\begin{equation}
|\Psi^a_i(k^+,z')\rangle_{gg\rho}=\frac{ig}{2\pi^{3/2}}\int_{\Lambda}^{\vee}\frac{dp^+}{\sqrt{p^+}}\int_{x,z}\KK^j(x-z)\,\rho^b(x)\,|g^a_i(k^+,z')\,g^b_j(p^+,z)\rangle .
\label{eq:ggrho1}
\end{equation}

\paragraph*{Normalization.} The normalization of the dressed gluon to order $g^2$ gives the same $\mf N$ as Eq.~\eqref{eq:NX}, as it must, since one operator $\Omega$ dresses all soft states. The norm of the $\rho$-independent two-gluon component~\eqref{eq:gg} is proportional to $\int d\xi\,\Phat(\xi)\int d^2\tilde{\mathbf p}/\tilde{\mathbf p}^2=\big[2\log\frac{k^+}{\Lambda}-\frac{11}{6}\big]\int d^2\tilde{\mathbf p}/\tilde{\mathbf p}^2$. Here $\Phat$ is the splitting function~\eqref{eq:Phat} and $\tilde{\mathbf p}$ is the relative transverse momentum of the pair. The transverse integral is scaleless and vanishes in dimensional regularization.
